\documentclass[12pt,a4paper]{article}

% Compile with: lualatex sheafification_sieve.tex
% or xelatex sheafification_sieve.tex
% Requires a system CJK font (Noto Serif CJK SC, Source Han Serif SC, etc.)

\usepackage{fontspec}
\usepackage{amsmath,amssymb,mathtools}
\usepackage{geometry}
\usepackage{enumitem}
\usepackage{microtype}

\geometry{margin=2.5cm}

\IfFontExistsTF{Noto Serif CJK SC}{
  \setmainfont{Noto Serif CJK SC}
}{
  \IfFontExistsTF{Source Han Serif SC}{
    \setmainfont{Source Han Serif SC}
  }{
    \IfFontExistsTF{Droid Sans Fallback}{
      \setmainfont{Droid Sans Fallback}
    }{
      \PackageWarning{sheaf}{No suitable CJK font found. Chinese may not render correctly.}
    }
  }
}

\title{用 sieve 构造 sheafification}
\author{}
\date{}

\begin{document}

\maketitle

设 \((\mathcal{C},J)\) 是一个 site，其中 \(J\) 是定义在 \(\mathcal{C}\) 上的 Grothendieck topology。我们完全用 sieve 来构造 sheafification。

\section{\(J\)-覆盖 sieve}

对 \(U\in\mathcal{C}\)，记
\[
yU=\mathcal{C}(-,U)
\]
为 representable presheaf。\(U\) 上的一个 sieve 是 \(yU\) 的子预层：
\[
S\hookrightarrow yU.
\]
等价地，对每个对象 \(V\)，有一个集合
\[
S(V)\subseteq \mathcal{C}(V,U),
\]
并且满足拉回封闭性：
\[
f\in S(V),\quad g:W\to V
\quad\Longrightarrow\quad
f\circ g\in S(W).
\]
Grothendieck topology \(J\) 给每个 \(U\) 指定一族覆盖 sieve：
\[
J(U)\subseteq \operatorname{Sieve}(U).
\]
我们写
\[
\operatorname{Cov}_J(U):=J(U).
\]
这里完全不需要引入覆盖态射族。

\section{sieve 上的 matching family}

设
\[
F:\mathcal{C}^{op}\to\mathbf{Set}
\]
是一个 presheaf，\(S\) 是 \(U\) 上的 sieve。

一个定义在 \(S\) 上的 \(F\)-matching family，就是一个自然变换
\[
\alpha:S\to F.
\]
把它逐点写出来，就是对每个
\[
f:V\to U\in S(V)
\]
指定一个元素
\[
\alpha_f\in F(V),
\]
并且满足对任意 \(g:W\to V\)：
\[
\alpha_{f\circ g}=F(g)(\alpha_f).
\]
因此所有 matching families 构成集合
\[
\operatorname{Match}_F(S)
:=
\operatorname{Nat}(S,F).
\]
注意，这里 \(S\) 本身已经包含了所有复合箭头，所以相容性完全由自然变换条件表达。

\section{对覆盖 sieve 取 filtered colimit}

若 \(S,T\in J(U)\)，且
\[
S\subseteq T,
\]
则有 restriction map
\[
\operatorname{Nat}(T,F)
\longrightarrow
\operatorname{Nat}(S,F),
\qquad
\alpha\longmapsto\alpha|_S.
\]
所以得到一个函子
\[
J(U)^{op}\longrightarrow\mathbf{Set},
\qquad
S\longmapsto\operatorname{Nat}(S,F).
\]
定义第一次 plus construction：
\[
F^+(U)
:=
\varinjlim_{S\in J(U)^{op}}
\operatorname{Nat}(S,F).
\]
为什么使用 \(J(U)^{op}\)？因为 sieve 的包含关系给出的是限制：
\[
S\subseteq T
\quad\Rightarrow\quad
\operatorname{Nat}(T,F)\to\operatorname{Nat}(S,F).
\]

\section{这个 colimit 的具体等价关系}

\(F^+(U)\) 的元素可以写成等价类
\[
[S,\alpha],
\]
其中
\[
S\in J(U),\qquad \alpha:S\to F.
\]
两个代表
\[
[S,\alpha],
\qquad
[T,\beta]
\]
相等，当且仅当它们在某个共同的覆盖 refinement 上相等。

具体地说，存在
\[
R\in J(U)
\]
满足
\[
R\subseteq S\cap T
\]
并且
\[
\alpha|_R=\beta|_R.
\]
这里 \(S\cap T\) 仍然是覆盖 sieve，这是 Grothendieck topology 的有限交性质：
\[
S,T\in J(U)
\quad\Longrightarrow\quad
S\cap T\in J(U).
\]
所以 \(J(U)^{op}\) 是 filtered 的：任意两个覆盖 sieve \(S,T\) 在反向包含关系下都有共同上界 \(S\cap T\)。

因此 \(F^+(U)\) 的含义是：
\begin{quote}
在某个覆盖 sieve 上定义的 matching family，模掉在更细覆盖 sieve 上相等的关系。
\end{quote}

\section{\(F^+\) 的 presheaf 结构}

给定箭头
\[
u:V\to U
\]
和 \(U\) 上的 sieve \(S\)，定义其拉回 sieve
\[
u^*S
\]
为
\[
(u^*S)(W)
=
\{g:W\to V\mid u\circ g\in S(W)\}.
\]
Grothendieck topology 的拉回公理保证：
\[
S\in J(U)
\quad\Longrightarrow\quad
u^*S\in J(V).
\]
若 \(\alpha:S\to F\)，定义
\[
u^*\alpha:u^*S\to F
\]
为
\[
(u^*\alpha)_g
=
\alpha_{u\circ g}.
\]
于是定义
\[
F^+(u):F^+(U)\to F^+(V)
\]
为
\[
[S,\alpha]
\longmapsto
[u^*S,u^*\alpha].
\]
因此 \(F^+\) 是一个 presheaf。

\section{自然映射 \(F\to F^+\)}

对 \(x\in F(U)\)，可以在最大 sieve
\[
yU=\mathcal{C}(-,U)
\]
上定义 matching family
\[
\alpha^x:yU\to F
\]
如下：
\[
\alpha^x_f=F(f)(x),
\qquad f:V\to U.
\]
于是有自然映射
\[
\eta_F:F\to F^+
\]
定义为
\[
\eta_{F,U}(x)
=
[yU,\alpha^x].
\]
这一步把一个真正定义在 \(U\) 上的截面，视为最大 sieve 上的 matching family。

\section{第一次 plus 后得到 separated presheaf}

一个 presheaf \(P\) 称为 \(J\)-separated，如果对每个覆盖 sieve \(S\in J(U)\)，限制映射
\[
P(U)\longrightarrow \operatorname{Nat}(S,P)
\]
都是单射。

也就是说，如果 \(x,y\in P(U)\) 在覆盖 sieve \(S\) 上的限制相同，那么
\[
x=y.
\]
重要事实是：
\[
F^+\text{ 总是 }J\text{-separated}.
\]
直观原因是，\(F^+(U)\) 中的元素本来就是“局部 matching family 的等价类”。如果两个元素在一个覆盖 sieve 上相等，就把它们代表的 sieve、这个覆盖 sieve，以及它们相等的部分取交，得到一个共同的覆盖 refinement；于是两个等价类本身相等。

但是 \(F^+\) 一般不一定已经是 sheaf。它具有唯一性，但不一定具有存在性。

\section{第二次 plus：sheafification}

再次应用同一个构造：
\[
aF:=(F^+)^+.
\]
即对每个 \(U\)：
\[
aF(U)
=
\varinjlim_{S\in J(U)^{op}}
\operatorname{Nat}(S,F^+).
\]
因此 \(aF(U)\) 的元素是：
\[
[S,\alpha],
\]
其中
\[
S\in J(U),\qquad
\alpha:S\to F^+.
\]
也就是说，对每个
\[
f:V\to U\in S(V),
\]
\(\alpha\) 给出一个元素
\[
\alpha_f\in F^+(V),
\]
并且这些元素满足 sieve 意义下的自然性条件：
\[
\alpha_{f\circ g}=F^+(g)(\alpha_f).
\]
最后再模掉在某个更细覆盖 sieve 上相等的 matching families。

结论是：
\[
\boxed{aF=(F^+)^+}
\]
是一个 \(J\)-sheaf。

\section{为什么第二次 plus 后是 sheaf}

第一次 plus 得到的 \(F^+\) 已经是 separated 的。

现在取 \(U\) 上的覆盖 sieve \(S\in J(U)\)，以及一个 matching family
\[
\alpha:S\to (F^+)^+.
\]
对每个
\[
f:V\to U\in S(V),
\]
\(\alpha_f\) 是 \((F^+)^+(V)\) 中的元素，因此它可以由某个 \(V\) 上的覆盖 sieve
\[
R_f\in J(V)
\]
以及 matching family
\[
\beta_f:R_f\to F^+
\]
表示。

利用：
\begin{enumerate}[label=\arabic*.]
\item \(J\) 的拉回稳定性；
\item \(J\) 的传递性；
\item \(F^+\) 的 separated 性；
\end{enumerate}
可以把这些局部代表放到一个共同的覆盖 sieve 上，并且把它们拼成一个 \(F^+\)-值 matching family。这个 matching family 给出一个元素
\[
x\in (F^+)^+(U).
\]
它限制到 \(S\) 后正好等于原来的 \(\alpha\)，因此存在性成立。

唯一性来自 \(F^+\) 的 separated 性：两个可能的粘合如果在 \(S\) 上相等，那么它们在某个覆盖 refinement 上相等，因而在 colimit 中代表同一个元素。

所以 \((F^+)^+\) 同时满足：
\begin{itemize}
\item matching family 存在粘合；
\item 粘合唯一。
\end{itemize}
因此它是 sheaf。

\section{Sheafification 的泛性质}

自然映射
\[
F\longrightarrow F^+\longrightarrow (F^+)^+
\]
给出
\[
\eta:F\to aF.
\]
对任意 \(J\)-sheaf \(G\)，复合 \(\eta\) 给出双射
\[
\operatorname{Hom}_{\operatorname{Sh}(\mathcal{C},J)}(aF,G)
\cong
\operatorname{Hom}_{\operatorname{PSh}(\mathcal{C})}(F,G).
\]
因此 \(aF\) 是 \(F\) 到 sheaf 范畴中的反射：
\[
\operatorname{PSh}(\mathcal{C})
\;\underset{\text{inclusion}}{\overset{a}{\rightleftarrows}}\;
\operatorname{Sh}(\mathcal{C},J).
\]
完全用 sieve 表述，sheafification 的核心公式就是：
\[
\boxed{
F^+(U)
=
\varinjlim_{S\in J(U)^{op}}
\operatorname{Nat}(S,F)
}
\]
以及
\[
\boxed{
aF=(F^+)^+
}.
\]
其中：
\begin{itemize}
\item \(S\in J(U)\) 是 \(U\) 上的覆盖 sieve；
\item \(\operatorname{Nat}(S,F)\) 是 \(S\) 上的 \(F\)-值 matching families；
\item \(J(U)^{op}\) 是 filtered category；
\item colimit 把在更细覆盖 sieve 上相等的 matching families 识别；
\item 第一次 plus 强制唯一性；
\item 第二次 plus 补上存在性；
\item 最终得到 sheafification。
\end{itemize}

\end{document}